\section{Differential properties of smooth morphisms}
\label{II.4}

To simplify, we shall restrict ourselves for the most part to
differential calculus of order~$1$, confining ourselves to brief
indications for higher order (where the results are just as simple).

For the definition of the sheaf of $1$-differentials
$\it{\Omega}^1_{X/Y}$ of a $Y$-prescheme~$X$,
\cf (I~\No\SourceRef{I.1}{1}). Suppose that $X$ and $Y$ are
$S$-preschemes, the structure morphism $f\colon X\to Y$
being an $S$-morphism. Then $f$ defines a homomorphism of
Modules (compatible with~$f$)
\begin{equation}
\label{eq:II.4.1}
f^*\colon \it{\Omega}^1_{Y/S}\to\it{\Omega}^1_{X/S}
\end{equation}
in other words, $\it{\Omega}^1_{X/S}$ is \emph{contravariant} in
the $S$-prescheme~$X$. Moreover~\eqref{eq:II.4.1}
is equivalent to a homomorphism of Modules on~$X$
\begin{equation*}
\label{eq:II.4.1bis}
\tag{4.1\textrm{bis}}
f^*\big(\it{\Omega}^1_{Y/S}\big)\to\it{\Omega}^1_{X/S}
\end{equation*}
also denoted by~$f^*$ for want of a better notation, and
which fits into a canonical exact sequence of homomorphisms of
Modules
\begin{equation}
\label{eq:II.4.2}
f^*\big(\it{\Omega}^1_{Y/S}\big)\to
\it{\Omega}^1_{X/S}\to\it{\Omega}^1_{X/Y}\to 0
\end{equation}
All these homomorphisms are defined by the condition of being of
local nature (which reduces one to the affine case) and of commuting
with the operators~$\rd$. The exactness of~\eqref{eq:II.4.2} is
classical and trivial, and transcribes in the affine case as the
exact sequence (corresponding to a homomorphism $B\to C$ of
$A$-algebras):
\begin{equation*}
\label{eq:II.4.2bis}
\tag{4.2\textrm{bis}}
\Omega^1_{B/A}\otimes_{B}C\to\Omega^1_{C/A}\to\Omega^1_{C/B}\to 0
\end{equation*}

\begin{lemma}
\label{II.4.1}
Let
% original p. 36
$f\colon X\to Y$ be a morphism of $S$-preschemes. If $f$
is unramified (\resp \'etale) then
$f^*\big(\it{\Omega}^1_{Y/S}\big)\to\it{\Omega}^1_{X/S}$
is surjective (\resp an isomorphism). The converse is true
in the ``unramified'' case, if $f$ is assumed locally of
finite type.
\end{lemma}

The unramified case follows from the exact
sequence~\eqref{eq:II.4.2} and from (I~\SourceRef{I.3.1}{3.1}), but can also be
seen directly as the \'etale case. Consider the diagram
\begin{equation*}
\xymatrix@C=1.5cm{
X \ar[r]^-{\Delta_{X/Y}} & X\times_{Y}X \ar[d] \ar[r] &
X\times_{S}X \ar[d] \\ & Y \ar[r]^-{\Delta_{Y/S}} & Y\times_{S}Y
}
\end{equation*}
in which $X\times_{Y}X$ is identified with the fibered product of~$Y$ and
$X\times_{S}X$ over~$Y\times_{S}Y$. Since $f$ is unramified,
$X\to X\times_{Y}X$ is an open immersion, hence the
``conormal'' sheaf of the composite immersion $\Delta_{X/S}$ of the
latter with $X\times_{Y}X\to X\times_{S}X$ is isomorphic to
the inverse image on~$X$ of the conormal sheaf for the immersion
$X\times_{Y}X\to X\times_{S}X$. On the other hand, $X\to Y$ being
\'etale hence flat, $X\times_{S}X\to Y\times_{S}Y$ is flat, hence the
conormal sheaf for the immersion $X\times_{Y}X\to X\times_{S}X$ is
isomorphic to the inverse image of the conormal sheaf for the immersion
$Y\to Y\times_{S}Y$, \ie the inverse image of~$\it{\Omega}^1_{Y/S}$. The
conclusion follows from this.

\begin{lemma}
\label{II.4.2}
Let $X=Y[t_{1},\dots,t_{n}]$, $Y$ being an
$S$-prescheme. Then the sequence of canonical homomorphisms
\[
0\to f^*\big(\it{\Omega}^1_{Y/S}\big)\to
\it{\Omega}^1_{X/S}\to\it{\Omega}^1_{X/Y}\to 0
\]
is exact and $\it{\Omega}^1_{X/Y}$ is free with basis the
$\rd_{X/Y}t_{i}$.
\end{lemma}

The verification (purely affine) is immediate. (N.B.\ one
already knows the exactness of~\eqref{eq:II.4.2}).

Combining these two statements and definition~\ref{II.1.1}, one
finds

\begin{theorem}
\label{II.4.3}
Let $f\colon X\to Y$ be a smooth morphism of $S$-preschemes,
then:
\begin{enumerate}
\item[(i)] The sequence of canonical homomorphisms
\[
0\to f^*\big(\it{\Omega}^1_{Y/S}\big)\to
\it{\Omega}^1_{X/S}\to\it{\Omega}^1_{X/Y}\to 0
\]
is exact.
\item[(ii)] $\it{\Omega}^1_{X/Y}$ is locally free, its rank $n$
at~$x$ is equal to the relative dimension of~$f$ at~$x$.
\end{enumerate}
\end{theorem}

\begin{corollary}
\label{II.4.4}
The homomorphism
% original p. 37
$f^*\big(\it{\Omega}^1_{Y/S}\big)\to\it{\Omega}^1_{X/S}$ is
injective, its image in $\it{\Omega}^1_{X/S}$ is locally a
direct factor.
\end{corollary}

Let $u\colon F\to G$ be a homomorphism of Modules on the
prescheme~$X$; one says that it is \emph{universally
injective}
at $x\in X$ if the homomorphism $F_{x}\to G_{x}$ of
$\cal{O}_{x}$-modules is injective, and remains so upon tensoring
with any $\cal{O}_{x}$-algebra (or, which comes to the same thing
moreover, with any $\cal{O}_{x}$-module). It suffices for example
that there exist an open neighborhood $U$ of~$x$ such that $u$ induces an
isomorphism of~$F|U$ onto a direct factor of~$G|U$; this condition
is also necessary when $F$ and $G$ are free (and of finite
type) in a neighborhood of~$x$; more precisely in this case
the following conditions are equivalent:
\begin{enumerate}
\item[(i)] $u$ is injective at~$x$ and $\Coker u$ free at~$x$;
\item[(ii)] There exists an open neighborhood $U$ of~$x$ such that $u$
induces an isomorphism of~$F|U$ onto a direct factor of~$G|U$;
\item[(iii)] $u$ is universally injective at~$x$;
\item[(iv)] the homomorphism
$F_{x}\otimes \kres(x)\to G_{x}\otimes \kres(x)$
on the restricted ``fibers'' induced by~$u$ is injective;
\item[(v)] The transposed homomorphism $\check{G}\to\check{F}$ is
surjective at the point~$x$ (or again, which comes to the
same thing, in a neighborhood of~$x$).
\end{enumerate}
(Circular proof, (iv)$\To$(v) follows from
Nakayama, on the other hand (v)$\To$(i)
since a locally free quotient sheaf is necessarily a
direct factor). Geometrically, the situation envisaged
means that $u$ corresponds to an isomorphism of the vector bundle whose
sheaf of sections is~$F$, onto a subbundle of the
analogous vector bundle defined by~$G$. Of course, it does
not suffice for this that $F\to G$ be injective.

\begin{corollary}
\label{II.4.5}
Let $f\colon X\to Y$ be a morphism of $S$-preschemes,
locally of finite type, $x\in X$, $y=f(x)$, $s$ the projection of~$x$
and~$y$ on~$S$. One assumes $Y$ is smooth at~$y$ over~$S$. Equivalent
conditions:
\begin{enumerate}
\item[(i)] $f$ is smooth at~$x$.
\item[(ii)] $X$ is smooth over $S$
at~$x$, and
$f^*\big(\it{\Omega}^1_{Y/S}\big)\to\it{\Omega}^1_{X/S}$ is
universally injective at~$x$, \ie it is an injective homomorphism at~$x$ and its cokernel $\it{\Omega}^1_{X/Y}$ is free
at~$x$.
\end{enumerate}
\end{corollary}

Necessity follows from~\ref{II.1.3}~(iii) and
from~\ref{II.4.3}~(i)\,(ii); let us prove sufficiency. Since the
$\rd g$ ($g\in\cal{O}_{x}$) generate the module
$\it{\Omega}^1_{X/Y}$ at~$x$, one can find
% original p. 38
$g_{i}$ ($1\le i\le n$) such that the images of the $\rd g_{i}$ in
$\big(\it{\Omega}^1_{X/Y}\big)_{x}$ form a basis of this
module. Taking $X$ small enough, one may assume that the $g_{i}$
come from sections of~$\cal{O}_{X}$, and thus define a
$Y$-morphism $g\colon X\to Y'=Y[t_{1},\dots,t_{n}]$. Using
the hypothesis and lemma~\ref{II.4.2}, one sees easily that
the corresponding homomorphism
$g^*\big(\it{\Omega}^1_{Y'/S}\big)\to\it{\Omega}^1_{X/S}$ is
bijective at~$x$, which reduces us to proving the

\begin{corollary}
\label{II.4.6}
Let $f\colon X\to Y$ be a morphism of smooth $S$-preschemes.
For $f$ to be \'etale at $x\in X$, it is necessary and sufficient
that $f^*\big(\it{\Omega}^1_{Y/S}\big)\to\it{\Omega}^1_{X/S}$ be
an isomorphism at~$x$.
\end{corollary}

One knows that this is necessary by~\ref{II.4.1}, and this condition
implies that $f$ is unramified at~$x$ by the same lemma. By virtue
of~\ref{II.2.2}, one is reduced to the case where
$S=\Spec(k)$. Since $Y$ is smooth over~$k$, it is regular, hence a
fortiori normal, and by virtue of (I~\SourceRef{I.9.5}{9.5}~(ii)) one is reduced
to proving that $\cal{O}_{y}\to\cal{O}_{x}$ is injective, or equivalently
that $\cal{O}_{y}$ and $\cal{O}_{x}$ have the same dimension. Now these
dimensions are respectively the ranks of~$\it{\Omega}^1_{Y/k}$ and
$\it{\Omega}^1_{X/k}$ at~$y$ \resp $x$, hence equal by virtue of
the hypothesis.
\begin{remarks}
\label{II.4.7}
$X$ and $Y$ being assumed smooth over~$S$, the
criterion~\ref{II.4.5}~(ii) of smoothness of $f\colon X\to Y$ can
also be stated by saying that for every $x\in X$, the
\emph{tangent} map (relative to the base~$S$) of~$f$ at~$x$,
\ie the transpose of the homomorphism of the $\kres(x)$ finite-dimensional vector spaces,
restricted fibers of
$f^*\big(\it{\Omega}^1_{Y/S}\big)$ and $\it{\Omega}^1_{X/S}$ at
$x$, is \emph{surjective}. This is a quite
familiar hypothesis in particular among people working with
analytic spaces. The non-singularity hypothesis which they
ordinarily make (which means that $X$ and $Y$ are ``smooth over~$\CC$'', \cf \No\ref{II.5}) seems due only to the fear which
singular points of algebraic varieties or analytic spaces still inspire
in many geometers.
\end{remarks}

Let us point out the following particular case of~\ref{II.4.6}:
\begin{corollary}
\label{II.4.8}
Let $X$ be an $S$-prescheme, $g\colon X\to
S[t_{1},\dots,t_{n}]$ an $S$-morphism, defined by the sections
$g_{i}$ ($1\le i\le n$) of~$\cal{O}_{X}$, $x$ a point of~$X$ such that
$X$ is smooth over~$S$ at~$x$. For $g$ to be \'etale at~$x$, it
is necessary and sufficient that the $\rd g_{i}$
% original p. 39
($1\le i\le n$) form a basis of~$\it{\Omega}^1_{X/S}$ at~$x$ (or,
which comes to the same thing, that their images in
$\it{\Omega}^1_{X/S}(x) =
\big(\it{\Omega}^1_{X/S}\big)_{x}\otimes_{\cal{O}_{x}}\kres(x)$ form
a basis of this vector space over~$\kres(x)$).
\end{corollary}

Let $X$ be a prescheme, $Y$ a closed subprescheme
of~$X$ defined by a coherent sheaf $\cal{J}$
of ideals. Thus $\cal{J}/\cal{J}^2$ may be regarded
as a coherent sheaf on~$Y$ (the \emph{conormal sheaf}
of~$Y$ in~$X$). If now $X$ is an $S$-prescheme, one has
a canonical exact sequence of quasi-coherent sheaves on~$Y$
\begin{equation}
\label{eq:II.4.3}
\cal{J}/\cal{J}^2\lto{\rd}
\it{\Omega}^1_{X/S}\otimes_{\cal{O}_{X}}\cal{O}_{Y}\to
\it{\Omega}^1_{Y/S}\to 0
\end{equation}
whose right-hand part is none other than~\eqref{eq:II.4.2} (with the
role of~$X$ and~$Y$ interchanged, taking into account that
$\it{\Omega}^1_{Y/X}=0$), while the homomorphism
$\cal{J}/\cal{J}^2\to\it{\Omega}^1_{X/S}\otimes_{\cal{O}_{X}}\cal{O}_{Y}$
is deduced from the homomorphism (in general not
linear) $g\to\rd g$ by passing to
quotients. The exactness of~\eqref{eq:II.4.3} is classical and
moreover trivial, and is interpreted in the affine case by the
following exact sequence (corresponding to a surjective homomorphism
$B\to C$ of $A$-algebras, with kernel~$J$):
\begin{equation*}
\label{eq:II.4.3bis}
\tag{4.3\textrm{bis}}
J/J^2\to\Omega^1_{B/A}\otimes_{B}C\to\Omega^1_{C/A}\to 0\qquad
(C=B/J)
\end{equation*}
(an exact sequence which had already been used
implicitly in the proof of (I~\SourceRef{I.7.2}{7.2})!).
